\PassOptionsToPackage{table}{xcolor}
\documentclass[10pt,aspectratio=169]{beamer}
\newcommand{\LectureNumber}{04}
\input{course-theme.tex}
\title{Variables and console input}
\subtitle{CSC103 Programming Fundamentals / Lecture 04}
\author{Dr. Muhammad Farhan}
\institute{Associate Professor of Computer Science\\Department of Computer Science\\COMSATS University Islamabad\\Sahiwal Campus}
\date{Fall 2026}
\begin{document}
\begin{frame}\titlepage\end{frame}

\begin{frame}[fragile,t]{The problem for this lecture}
\begin{columns}[T,onlytextwidth]\begin{column}{0.60\textwidth}
Read a quantity and price using Scanner. Print their product. Then read a one-letter category. Use quantity 3, price 12.5, category A.
\end{column}\begin{column}{0.35\textwidth}\small
Create one Scanner and read quantity, price and category in order.
\end{column}\end{columns}
\note{Introduce the target problem before explaining the tools. Revisit it in the guided solution later.\par Syllabus reading: Liang Ch 2. CLO-1.}
\end{frame}

\begin{frame}[fragile,t]{Learning objectives}
\begin{columns}[T,onlytextwidth]\begin{column}{0.60\textwidth}
\begin{itemize}
\item Declare and initialise variables and constants
\item Read numeric and character input
\item Trace assignment as a change of stored value
\end{itemize}
\end{column}\begin{column}{0.35\textwidth}\small
CLO-1

Variables and initialisation\par Assignment copies a value\par Scanner input\par Reading one character
\end{column}\end{columns}
\note{Explain how each objective supports the opening problem. The final questions check these objectives.\par Syllabus reading: Liang Ch 2. CLO-1.}
\end{frame}

\begin{frame}[fragile,t]{Variables and initialisation}
\begin{itemize}
\item A variable has a declared type and a name.
\item Initialisation supplies its first value. A local variable must receive a value before a read.
\end{itemize}
\vspace{1mm}\begin{block}{Example}
\begin{lstlisting}
int seats = 30;
seats = 28;
final int ROOM_LIMIT = 40;
\end{lstlisting}
\end{block}
\note{A final variable can be assigned once. Explain the distinction between a variable value changing and a constant preserving a fixed value.\par Syllabus reading: Liang Ch 2. CLO-1.}
\end{frame}

\begin{frame}[fragile,t]{Assignment copies a value}
\begin{itemize}
\item The right-hand expression is evaluated before assignment.
\item For primitive variables, assigning one variable to another copies the current primitive value.
\end{itemize}
\vspace{1mm}\begin{block}{Example}
\begin{lstlisting}
int a = 7;
int b = a;
a = 9;
// b is still 7
\end{lstlisting}
\end{block}
\note{Use a board trace with a and b as separate named cells. This is a conceptual model, not a claim about exact physical memory layout.\par Syllabus reading: Liang Ch 2. CLO-1.}
\end{frame}

\begin{frame}[fragile,t]{Scanner input}
\begin{itemize}
\item Create one Scanner for console input.
\item nextInt reads an integer token and nextDouble reads a decimal token.
\item The entered token must match the requested type.
\end{itemize}
\vspace{1mm}\begin{block}{Example}
\begin{lstlisting}
java.util.Scanner in =
    new java.util.Scanner(System.in);
int age = in.nextInt();
\end{lstlisting}
\end{block}
\note{Give an import java.util.Scanner alternative. State that exception handling for malformed input appears later. Do not create several wrappers around System.in.\par Syllabus reading: Liang Ch 2. CLO-1.}
\end{frame}

\begin{frame}[fragile,t]{Reading one character}
\begin{itemize}
\item Scanner has no nextChar method.
\item next().charAt(0) reads the first UTF-16 code unit of the next token.
\item For this exercise, input is one basic Latin character.
\end{itemize}
\vspace{1mm}\begin{block}{Example}
\begin{lstlisting}
String token = in.next();
char choice = token.charAt(0);
\end{lstlisting}
\end{block}
\note{The token cannot be empty after a successful next(). A multi-character token gives only its first char. A Java char does not necessarily represent an entire Unicode character.\par Syllabus reading: Liang Ch 2. CLO-1.}
\end{frame}

\begin{frame}[fragile,t]{Choosing a type from the data}
\begin{itemize}
\item A quantity that counts whole items fits an integer type.
\item A price with a fractional part needs a suitable numeric representation.
\item For introductory arithmetic, double is convenient, but exact financial applications need decimal handling.
\end{itemize}
\vspace{1mm}\begin{block}{Example}
\begin{lstlisting}
quantity: 3
price: 12.5
category: A
Each input has a different role and type.
\end{lstlisting}
\end{block}
\note{Explain the mechanism using the displayed example. A quantity that counts whole items fits an integer type. A price with a fractional part needs a suitable numeric representation. For introductory arithmetic, double is convenient, but exact financial applications need decimal handling.\par Syllabus reading: Liang Ch 2. CLO-1.}
\end{frame}

\begin{frame}[fragile,t]{Console input is a sequence of tokens}
\begin{itemize}
\item Scanner reads the next available token when a next method succeeds.
\item The program must request the token type in the order the user supplies it.
\item A character read through next().charAt(0) uses only the first character of that token.
\end{itemize}
\vspace{1mm}\begin{block}{Example}
\begin{lstlisting}
Input tokens: 3 12.5 A
nextInt consumes 3
nextDouble consumes 12.5
next consumes A
\end{lstlisting}
\end{block}
\note{Explain the mechanism using the displayed example. Scanner reads the next available token when a next method succeeds. The program must request the token type in the order the user supplies it. A character read through next().charAt(0) uses only the first character of that token.\par Syllabus reading: Liang Ch 2. CLO-1.}
\end{frame}


\begin{frame}[fragile,t]{Input tokens become typed values}
\begin{center}\begin{tikzpicture}
\node[box] (n0) at (0.0,0) {12  2.5  A};
\node[box] (n1) at (3.45,0) {nextInt: 12};
\node[box] (n2) at (6.9,0) {nextDouble: 2.5};
\node[alt] (n3) at (10.350000000000001,0) {next: A};
\draw[arr] (n0) -- (n1);
\draw[arr] (n1) -- (n2);
\draw[arr] (n2) -- (n3);
\end{tikzpicture}\end{center}
\begin{block}{What to notice}Each Scanner call consumes the next token according to the requested type.\end{block}
\note{Trace each labelled connection. Each Scanner call consumes the next token according to the requested type.}
\end{frame}
\begin{frame}[fragile,t]{Key distinctions}
\small\renewcommand{\arraystretch}{1.5}
\rowcolors{2}{pale}{white}
\begin{tabularx}{\textwidth}{>{\raggedright\arraybackslash}X>{\raggedright\arraybackslash}X>{\raggedright\arraybackslash}X}
\rowcolor{navy}
\color{white}\bfseries Declaration & \color{white}\bfseries Initial value & \color{white}\bfseries May be reassigned?\\
int quantity = 3 & 3 & Yes\\
double price = 12.5 & 12.5 & Yes\\
char category = 'A' & A & Yes\\
final int LIMIT = 20 & 20 & No after initial assignment\\
\end{tabularx}
\vspace{3mm}\footnotesize These distinctions determine which operation or control structure fits the problem.
\note{\par Syllabus reading: Liang Ch 2. CLO-1.}
\end{frame}

\begin{frame}[fragile,t]{First example: execution}
\begin{lstlisting}
int initial = 10;
int remaining = initial;
remaining = remaining - 3;
final int CAPACITY = 20;
System.out.println(remaining);
System.out.println(initial);
\end{lstlisting}
\note{This is the main body from Java\_Examples/Lecture04.java. The source file includes the complete class. Predict the result before the next slide.\par Syllabus reading: Liang Ch 2. CLO-1.}
\end{frame}

\begin{frame}[fragile,t]{First example: result}
\begin{columns}[T,onlytextwidth]\begin{column}{0.60\textwidth}
7\par 10\par Assigning remaining does not change initial.
\end{column}\begin{column}{0.35\textwidth}\small
Scanner input

Reading one character
\end{column}\end{columns}
\note{Expected output and interpretation: 7
10
Assigning remaining does not change initial.\par Syllabus reading: Liang Ch 2. CLO-1.}
\end{frame}

\begin{frame}[fragile,t]{Guided practice}
\begin{columns}[T,onlytextwidth]\begin{column}{0.60\textwidth}
Read a quantity and price using Scanner. Print their product. Then read a one-letter category. Use quantity 3, price 12.5, category A.
\end{column}\begin{column}{0.35\textwidth}\small
Attempt the solution before continuing.

Use the definitions and examples from this lecture.
\end{column}\end{columns}
\note{Give students 8 minutes to attempt the problem. The following slides provide a complete solution. Original solution guidance: Use int quantity = in.nextInt(), double price = in.nextDouble(), and char category = in.next().charAt(0). Product: 37.5. Category: A.\par Syllabus reading: Liang Ch 2. CLO-1.}
\end{frame}

\begin{frame}[fragile,t]{Solution strategy}
\begin{columns}[T,onlytextwidth]\begin{column}{0.60\textwidth}
\begin{itemize}
\item Create one Scanner and read quantity, price and category in order.
\item Calculate quantity * price using the numeric values.
\item Print the product and the selected category.
\end{itemize}
\end{column}\begin{column}{0.35\textwidth}\small
Sample console input\par 3 12.5 A\par 
\end{column}\end{columns}
\note{Discuss why each step is needed. State the input assumptions before executing the program.\par Syllabus reading: Liang Ch 2. CLO-1.}
\end{frame}

\begin{frame}[fragile,t]{Complete solution}
\begin{lstlisting}
public class Solution04 {
  public static void main(String[] args) throws Exception {
    java.util.Scanner in = new java.util.Scanner(System.in);
    in.useLocale(java.util.Locale.US);
    int quantity = in.nextInt();
    double price = in.nextDouble();
    char category = in.next().charAt(0);
    double total = quantity * price;
    System.out.println(total);
    System.out.println(category);
  }
}
\end{lstlisting}
\note{Complete runnable file: Java\_Examples/Solution04.java. Standard input: 3 12.5 A
  \par Syllabus reading: Liang Ch 2. CLO-1.}
\end{frame}

\begin{frame}[fragile,t]{Execution trace}
\small\renewcommand{\arraystretch}{1.5}
\rowcolors{2}{pale}{white}
\begin{tabularx}{\textwidth}{>{\raggedright\arraybackslash}X>{\raggedright\arraybackslash}X>{\raggedright\arraybackslash}X}
\rowcolor{navy}
\color{white}\bfseries Read or calculation & \color{white}\bfseries Consumed token & \color{white}\bfseries Stored result\\
nextInt() & 3 & quantity = 3\\
nextDouble() & 12.5 & price = 12.5\\
next().charAt(0) & A & category = 'A'\\
quantity * price & No new input & total = 37.5\\
\end{tabularx}
\vspace{3mm}\footnotesize Follow the changing state in execution order.
\note{Manually derived trace for the complete solution. Create one Scanner and read quantity, price and category in order. Calculate quantity * price using the numeric values. Print the product and the selected category.\par Syllabus reading: Liang Ch 2. CLO-1.}
\end{frame}

\begin{frame}[fragile,t]{Solution result}
\begin{columns}[T,onlytextwidth]\begin{column}{0.60\textwidth}
37.5\par A
\end{column}\begin{column}{0.35\textwidth}\small
Why this result follows

Print the product and the selected category.
\end{column}\end{columns}
\note{Expected result: 37.5
A\par Syllabus reading: Liang Ch 2. CLO-1.}
\end{frame}

\begin{frame}[fragile,t]{A common incorrect approach}
int total;\par System.out.println(total);
\vspace{1mm}\begin{block}{Example}
\begin{lstlisting}
A local variable has no usable automatic initial value. Assign total before printing it, for example int total = 0.
\end{lstlisting}
\end{block}
\note{Ask students to predict the failure before discussing the explanation. The right side gives the correction.\par Syllabus reading: Liang Ch 2. CLO-1.}
\end{frame}

\begin{frame}[fragile,t]{Boundary and failure checks}
\begin{columns}[T,onlytextwidth]\begin{column}{0.60\textwidth}
Check the stated input assumptions.

Read a rectangle length and width. Print area and perimeter. State the units and allowed input range.
\end{column}\begin{column}{0.35\textwidth}\small
Use int quantity = in.nextInt(), double price = in.nextDouble(), and char category = in.next().charAt(0). Product: 37.5. Category: A.
\end{column}\end{columns}
\note{Use the specification to derive each expected result. Exercise solution: Use int quantity = in.nextInt(), double price = in.nextDouble(), and char category = in.next().charAt(0). Product: 37.5. Category: A.\par Syllabus reading: Liang Ch 2. CLO-1.}
\end{frame}

\begin{frame}[fragile,t]{Check your understanding}
\begin{columns}[T,onlytextwidth]\begin{column}{0.60\textwidth}
After int x=4; int y=x; y=8;, what is x? Can a final variable be assigned again?
\end{column}\begin{column}{0.35\textwidth}\small
Write a prediction and explain the rule that supports it.
\end{column}\end{columns}
\note{Answer: x remains 4. A final variable cannot be reassigned after its initial assignment.\par Syllabus reading: Liang Ch 2. CLO-1.}
\end{frame}

\begin{frame}[fragile,t]{Answer and explanation}
\begin{columns}[T,onlytextwidth]\begin{column}{0.60\textwidth}
x remains 4. A final variable cannot be reassigned after its initial assignment.
\end{column}\begin{column}{0.35\textwidth}\small
A local variable has no usable automatic initial value. Assign total before printing it, for example int total = 0.
\end{column}\end{columns}
\note{Reconnect the answer to the relevant definition rather than treating it as a fact to memorise.\par Syllabus reading: Liang Ch 2. CLO-1.}
\end{frame}


\begin{frame}[fragile,t]{Compare cases and predict outcomes}
\small\renewcommand{\arraystretch}{1.5}\rowcolors{2}{pale}{white}
\begin{tabularx}{\textwidth}{>{\raggedright\arraybackslash}X>{\raggedright\arraybackslash}X>{\raggedright\arraybackslash}X}\rowcolor{navy}
\color{white}\bfseries Next token & \color{white}\bfseries Requested operation & \color{white}\bfseries Result\\
12 & nextInt() & int value 12\\
2.5 & nextDouble() & double value 2.5\\
A & next().charAt(0) & char value A\\
\end{tabularx}\par\vspace{3mm}\begin{exampleblock}{Discuss}Explain each expected result before running the example. Which case would reveal a wrong assumption?\end{exampleblock}
\note{Read the first two columns and invite predictions before discussing the outcome.}
\end{frame}

\begin{frame}[fragile,t]{Apply the idea}
\begin{alertblock}{Independent practice}A shop sells 12 items at 2.5 each. Show the variable declarations and total. Then predict what happens if nextInt() is used for the token 2.5.\end{alertblock}\vspace{3mm}\begin{enumerate}\item Work through the input by hand.\item Write or correct the program.\item Justify the expected result.\end{enumerate}\note{Allow 3–5 minutes, then compare answers in pairs. A shop sells 12 items at 2.5 each. Show the variable declarations and total. Then predict what happens if nextInt() is used for the token 2.5.}
\end{frame}

\begin{frame}[fragile,t]{Solution and reasoning}
\begin{lstlisting}
int quantity = 12;
double unitPrice = 2.5;
double total = quantity * unitPrice;
System.out.println(total);
\end{lstlisting}\vspace{1mm}\small\begin{itemize}
\item The int operand is promoted for multiplication with double.
\item The result is 30.0.
\item Scanner.nextInt() cannot consume 2.5 as an int and throws InputMismatchException.
\end{itemize}\note{Answer: The int operand is promoted for multiplication with double. The result is 30.0. Scanner.nextInt() cannot consume 2.5 as an int and throws InputMismatchException.}
\end{frame}

\begin{frame}[fragile,t]{Circle area with console input}
\begin{itemize}
\item Declaration introduces radius and area; assignment gives a local variable a usable value.
\item A named constant documents the chosen approximation to pi.
\item Read the radius as double, then calculate only after input is available.
\end{itemize}
\vspace{3mm}\footnotesize\textcolor{accent}{Source: supplied ISB campus slides. See speaker notes and ISB\_Source\_Map.md.}
\note{Adapted from the supplied ISB campus folder: Lecture{-}5.pptx, slides 3, 4, 5, 6, 12, 15, 22, 23. Adapted explanation and runnable implementation; sample inputs and traces added for this course.}
\end{frame}

\begin{frame}[fragile,t]{Worked problem: Circle area with console input}
\begin{block}{Task}Read radius 20 and calculate its area with PI = 3.14159, matching the source example. Assume a nonnegative radius.\end{block}
\small Input: 20\par\vspace{3mm}\begin{exampleblock}{Before running}Predict the output and identify the variables that change.\end{exampleblock}
\note{Adapted from the supplied ISB campus folder: Lecture{-}5.pptx, slides 3, 4, 5, 6, 12, 15, 22, 23. Adapted explanation and runnable implementation; sample inputs and traces added for this course.}
\end{frame}

\begin{frame}[fragile,t]{Complete ISB example (1/1)}
\begin{lstlisting}
public class IsbLecture04 {
  public static void main(String[] args) throws Exception {
    java.util.Scanner input = new java.util.Scanner(System.in);
    input.useLocale(java.util.Locale.US);
    final double PI = 3.14159;
    double radius = input.nextDouble();
    double area = PI * radius * radius;
    System.out.printf(java.util.Locale.US, "Area: %.3f%n", area);
  }

}
\end{lstlisting}
\note{Adapted from the supplied ISB campus folder: Lecture{-}5.pptx, slides 3, 4, 5, 6, 12, 15, 22, 23. Adapted explanation and runnable implementation; sample inputs and traces added for this course. Complete file: ISB\_Java\_Examples/IsbLecture04.java.}
\end{frame}

\begin{frame}[fragile,t]{Worked trace and expected result}
\small\renewcommand{\arraystretch}{1.4}\rowcolors{2}{pale}{white}
\begin{tabularx}{\textwidth}{>{\raggedright\arraybackslash}X>{\raggedright\arraybackslash}X>{\raggedright\arraybackslash}X}\rowcolor{navy}
\color{white}\bfseries Moment & \color{white}\bfseries radius & \color{white}\bfseries area\\
After reading & 20.0 & Not assigned yet\\
After calculation & 20.0 & 1256.636\\
Radius changed to zero & 0.0 & Recalculation gives 0.000\\
\end{tabularx}\par\vspace{2mm}
\begin{block}{Expected console output}\ttfamily\footnotesize Area: 1256.636\end{block}
\note{Adapted from the supplied ISB campus folder: Lecture{-}5.pptx, slides 3, 4, 5, 6, 12, 15, 22, 23. Adapted explanation and runnable implementation; sample inputs and traces added for this course. Trace and expected values independently worked for this edition.}
\end{frame}

\begin{frame}[fragile,t]{Extend the ISB example}
\begin{alertblock}{Practice}Why does changing radius after calculating area not automatically change area?\end{alertblock}\vspace{3mm}\begin{exampleblock}{Explain your reasoning}Assignment stores the computed value at that moment. Recompute area after changing radius.\end{exampleblock}
\note{Adapted from the supplied ISB campus folder: Lecture{-}5.pptx, slides 3, 4, 5, 6, 12, 15, 22, 23. Adapted explanation and runnable implementation; sample inputs and traces added for this course. Answer: Assignment stores the computed value at that moment. Recompute area after changing radius.}
\end{frame}
\begin{frame}[fragile,t]{Lecture summary}
\begin{columns}[T,onlytextwidth]\begin{column}{0.60\textwidth}
\begin{itemize}
\item and initialise variables and constants
\item numeric and character input
\item assignment as a change of stored value
\end{itemize}
\end{column}\begin{column}{0.35\textwidth}\small
\begin{itemize}
\item Create one Scanner and read quantity, price and category in order.
\item Calculate quantity * price using the numeric values.
\item Print the product and the selected category.
\end{itemize}
\end{column}\end{columns}
\note{Ask students to give one concrete example for the concepts listed. Use the worked solution to connect the topics.\par Syllabus reading: Liang Ch 2. CLO-1.}
\end{frame}

\begin{frame}[fragile,t]{After-class practice}
\begin{columns}[T,onlytextwidth]\begin{column}{0.60\textwidth}
Read a rectangle length and width. Print area and perimeter. State the units and allowed input range.
\end{column}\begin{column}{0.35\textwidth}\small
Syllabus reading\par Liang Ch 2

Next lecture: Types and arithmetic expressions
\end{column}\end{columns}
\note{Reading labels follow the supplied outline. Check topic headings against the available textbook edition. Require a program and expected results for the practice task.\par Syllabus reading: Liang Ch 2. CLO-1.}
\end{frame}
\end{document}
